\documentclass[addpoints,12pt,
noanswers]  % answers or noanswers
{exam}

\usepackage[utf8]{inputenc}
\usepackage[T1]{fontenc}
\usepackage{graphicx,amssymb,amsmath,multirow,tabularx,pgf,pgfpages}
\usepackage{libertine}
\usepackage[ngerman]{babel}

%%% Folgendes bitte anpassen %%%%%%%%%%%%%%%%%%%%%%%%%%%%%%%%%%%%%%%%
\newcommand{\theauthor}{Prof.\ Dr.\ Andreas Fischer}
\newcommand{\thedate}{10.02.2020}
\newcommand{\thetime}{11:00}
\newcommand{\semester}{SS 2019}
\newcommand{\course}{Echtzeitsysteme}
\newcommand{\programme}{BA Angewandte Informatik}
\newcommand{\examtype}{schr. P. 90 Min}
\newcommand{\duration}{90 Minuten}
\newcommand{\supplies}{Ein einseitig handschriftlich beschriebenes DINA4-Blatt}
\newcommand{\faculty}{Angewandte Informatik}
%%%%%%%%%%%%%%%%%%%%%%%%%%%%%%%%%%%%%%%%%%%%%%%%%%%%%%%%%%%%%%%%%%%%%%

\vqword{Frage}
\vpword{Punkte}
\vsword{Erzielt}
\vtword{Gesamt:}
\hqword{Frage:}
\hpword{Punkte:}
\hsword{Erzielt:}
\htword{Gesamt}
\pointpoints{Punkt}{Punkte}
\qformat{\bfseries Frage \thequestion~(\totalpoints~\points):\hfill}
\newcommand{\Seiten}{Seite \thepage\ von \numpages}

\pgfpagesdeclarelayout{boxed}
{%
  \edef\pgfpageoptionborder{0pt}
}
{%
  \pgfpagesphysicalpageoptions
  {%
    logical pages=1,%
  }
  \pgfpageslogicalpageoptions{1}
  {%
    border code=\pgfsetlinewidth{1pt}\pgfstroke,%
    border shrink=\pgfpageoptionborder,%
    resized width=.93\pgfphysicalwidth,%
    resized height=.95\pgfphysicalheight,%
    center=\pgfpoint{.5\pgfphysicalwidth}{.5\pgfphysicalheight}%
  }%
}
\pgfpagesuselayout{boxed}
\extrawidth{3cm}
\extraheadheight[-2cm]{2cm}
\extrafootheight{-1cm}

\newsavebox{\headbox}
\savebox{\headbox}{%
	\renewcommand{\arraystretch}{1.7}
	\parbox{\textwidth}{%
	\begin{tabularx}{\linewidth}{|X|X|X|X|}
	\hline
	\multicolumn{2}{|X|}{%
	\multirow{2}{*}{%
		\includegraphics[width=.43\textwidth]{THD_Logo.pdf}
	}
	}
	& \multicolumn{2}{l|}{Matrikelnummer:}\\
	\cline{3-4}
	\multicolumn{2}{|p{.45\textwidth}|}{}
	& 
	\multicolumn{2}{p{.45\textwidth}|}{Platzziffer:} \\
	\hline
	\textbf{Kurs:} & \course & \textbf{Semester:} & \semester \\
	\hline
	
	\textbf{Prüfer:} & \multicolumn{3}{p{.7\textwidth}|}{\theauthor} \\
	\hline
	\end{tabularx}
	}
}

\chead[]{\usebox\headbox}
\cfoot{\Seiten}



\begin{document}

\renewcommand{\arraystretch}{1.7}
\hspace{-2em}
\begin{tabularx}{\textwidth}{|X|X|X|X|}
	\hline
	\multicolumn{2}{|X|}{%
	\multirow{2}{*}{%
		\includegraphics[width=.43\textwidth]{THD_Logo.pdf}
	}
	}
	& \multicolumn{2}{l|}{Matrikelnummer:}\\
	\cline{3-4}
	\multicolumn{2}{|p{.5\textwidth}|}{}
	& Platzziffer: & Punkte:\quad\quad Note:\quad\\
	\hline
	\multicolumn{4}{|c|}{\textbf{\faculty}}\\
	\hline
	\textbf{Kurs:} & \course & \textbf{Semester:} & \semester \\
	\hline
	\textbf{Studiengang:} & \programme & \textbf{Dauer/Anteil:} & \duration \\
	\hline
	\textbf{Prüfer:} & \theauthor & \textbf{Prüfungsdatum} & \thedate \\
	\hline
	\textbf{Hilfsmittel:} & \supplies & \textbf{Uhrzeit} & \thetime \\
	\hline
	\textbf{Prüfungsart:} & \examtype & \textbf{Anzahl d. Blätter} & \numpages \\
	\hline
\end{tabularx}

\vfill


\textbf{Bevor Sie beginnen:}
\begin{itemize}
	\item Kontrollieren Sie ob Ihr Angabenblatt vollständig ist.
	\item Schreiben Sie Ihre Platzziffer auf jede Seite.
	\item Lesen Sie sich alle Aufgaben durch und verschaffen Sie sich einen Überblick über den Umfang der Aufgaben.
	\item Stellen Sie sicher, dass Ihre Antworten vollständig und lesbar sind.
	\item Schreiben Sie nicht außerhalb des markierten Seitenrandes.
\end{itemize}

\textbf{Viel Erfolg!}

\vfill

\begin{center}
	\gradetable[h]{}
\end{center}

\vfill

\clearpage


\begin{questions}

	\question[5] Frage \dots \emph{Punkte müssen angegeben werden!}

	\question Frage mit Unterpunkten
	\begin{parts}
		\part[3] Etwas
		\part[2] und noch etwas
	\end{parts}
	\section*{Wichtiges Thema}
	\question[10] Wichtige Frage \dots

\end{questions}
\end{document}
